\documentclass[conference]{IEEEtran}

\usepackage[T1]{fontenc}
\usepackage[utf8]{inputenc}
\usepackage{lmodern}
\usepackage{cite}
\usepackage{microtype}
\usepackage{booktabs}
\usepackage{longtable}
\usepackage{array}
\usepackage{amsmath,amssymb}
\usepackage{graphicx}
\usepackage{xcolor}
\usepackage{hyperref}

\hypersetup{
  colorlinks=true,
  linkcolor=blue!50!black,
  citecolor=blue!50!black,
  urlcolor=blue!50!black,
<<<<<<< HEAD
  pdftitle={Recuperacion de Bundle a Producto en Moda},
  pdfauthor={Equipo Hack2026}
}

\title{Recuperacion de Bundle a Producto en Moda:\\Paper Tecnico y Documentacion del Sistema}
\author{
\IEEEauthorblockN{Equipo Hack2026}
\IEEEauthorblockA{Informe Tecnico Interno}
=======
  pdftitle={Bundle-to-Product Retrieval in Fashion},
  pdfauthor={Hack2026 Team}
}

\title{Bundle-to-Product Retrieval in Fashion:\\Technical Paper and System Documentation}
\author{
\IEEEauthorblockN{Hack2026 Team}
\IEEEauthorblockA{Internal Technical Report}
>>>>>>> e9d5d43845f1bcb17b1a4f5c33b124b3b6907bf4
}
\date{\today}

\begin{document}

\maketitle

\begin{abstract}
<<<<<<< HEAD
Este documento describe el sistema completo de retrieval de Hack2026 para descomposicion de bundles de moda, donde cada imagen de bundle debe asociarse con uno o mas identificadores de producto. La implementacion actual se apoya en OpenCLIP (Marqo Fashion SigLIP), con deteccion opcional de regiones de ropa y varias estrategias de inferencia (pipeline principal, per-crop v8, phase1 por seccion y top-1 por crop). Se detallan el preprocesamiento de datos, el aumento offline, el entrenamiento del modelo, la validacion local, la inferencia y la generacion de entregas. El paper se apoya en el codigo del repositorio, con enfasis en reproducibilidad y decisiones de ingenieria.
\end{abstract}

\begin{IEEEkeywords}
recuperacion visual, aprendizaje multimodal, OpenCLIP, IA en moda, metric learning, matching a nivel de objeto
=======
This document describes the full Hack2026 retrieval system for fashion bundle decomposition, where each bundle image must be matched to one or more product identifiers. The implementation combines a baseline pretrained visual retrieval pipeline and an advanced OpenCLIP training pipeline with optional clothing-region detection. We detail data preprocessing, offline augmentation, model training, local validation, inference, and submission generation. The paper is grounded on the repository codebase, with an emphasis on reproducibility and engineering decisions.
\end{abstract}

\begin{IEEEkeywords}
visual retrieval, multimodal learning, OpenCLIP, fashion AI, metric learning, object-level matching
>>>>>>> e9d5d43845f1bcb17b1a4f5c33b124b3b6907bf4
\end{IEEEkeywords}

\section{Introduccion}

En goofytex, la tarea se formula como retrieval visual multietiqueta: dada una imagen de bundle, el sistema debe recuperar y ordenar los \texttt{product\_asset\_id} mas probables, con un maximo de 15 predicciones por bundle.

El repositorio implementa un stack modular orientado a experimentacion rapida y despliegue reproducible:
\begin{enumerate}
  \item Preparacion de datos e imagenes (\texttt{download.sh}, \texttt{preprocess\_data.py}, \texttt{offline\_augment.py}).
  \item Entrenamiento de retrieval multimodal (\texttt{python -m src.train}, backend \texttt{src/models/retrieval\_openclip.py}).
  \item Inferencia y exportacion de submission (\texttt{src.infer}, \texttt{src.new\_infer}, \texttt{src.infer\_phase1}, \texttt{src.infer\_top1}, \texttt{src.infer\_openclip}).
\end{enumerate}

La documentacion adopta CRISP-DM como marco de trabajo para organizar decisiones tecnicas, trazabilidad y fases de iteracion. El objetivo de este paper es consolidar la metodologia y las tecnologias usadas en el codigo, sin incluir resultados de rendimiento.

<<<<<<< HEAD
\section{Datos y Preprocesamiento}

\subsection{Datasets de Origen}

Las entradas principales son cuatro CSV. Su rol funcional se resume en la Tabla~\ref{tab:datasets}.
=======
\section{Data and Preprocessing}

\subsection{Source Datasets}

The core inputs are four CSV files. Their functional role in the project is summarized in Table~\ref{tab:datasets}.
>>>>>>> e9d5d43845f1bcb17b1a4f5c33b124b3b6907bf4

\begin{table}[h]
  \centering
  \small
  \begin{tabular}{p{0.36\linewidth} p{0.56\linewidth}}
    \toprule
<<<<<<< HEAD
    \textbf{Archivo} & \textbf{Proposito} \\
    \midrule
    \texttt{data/bundles\_dataset.csv} & Catalogo de bundles y URLs de imagen. \\
    \texttt{data/product\_dataset.csv} & Catalogo de productos, URLs y descripcion textual. \\
    \texttt{data/bundles\_product\_match\_train.csv} & Enlaces supervisados bundle-producto para entrenamiento. \\
    \texttt{data/bundles\_product\_match\_test.csv} & Plantilla de test para predecir \texttt{product\_asset\_id}. \\
    \bottomrule
  \end{tabular}
  \caption{Datasets usados por entrenamiento e inferencia.}
  \label{tab:datasets}
\end{table}

\subsection{Snapshot Local (2026-02-28)}

Inspeccion directa del workspace:
\begin{itemize}
  \item Bundles en catalogo: 2,331.
  \item Productos en catalogo: 27,688.
  \item Relaciones train: 6,493 (1,876 bundles unicos, 4,012 productos unicos).
  \item Filas test: 455 (sin IDs de bundle fuera del catalogo).
  \item Imagenes locales: 2,331 bundles y 27,688 productos.
\end{itemize}

\subsection{Pipeline de Preprocesamiento}

\texttt{preprocess\_data.py} implementa validacion de esquema, chequeos de integridad referencial, split train/val por \texttt{bundle\_asset\_id}, descarga opcional de imagenes y generacion de manifiestos. Artefactos principales:
\begin{itemize}
  \item \texttt{data/manifests/train\_manifest.jsonl}
  \item \texttt{data/manifests/val\_manifest.jsonl}
  \item \texttt{data/labels.json}, \texttt{data/label2idx.json}, \texttt{data/stats.json}
\end{itemize}

En el estado actual, \texttt{data/stats.json} reporta 86 etiquetas limpias y 6,493 relaciones validas tras limpieza.

\subsection{Aumento Offline}

\texttt{offline\_augment.py} genera vistas sinteticas para bundles y productos con semillas deterministas por asset+indice de augmentation. El script soporta manifests CSV o JSONL y guarda:
=======
    \textbf{File} & \textbf{Purpose} \\
    \midrule
    \texttt{data/bundles\_dataset.csv} & Bundle catalog and bundle image URLs. \\
    \texttt{data/product\_dataset.csv} & Product catalog, product image URLs, and textual descriptions. \\
    \texttt{data/bundles\_product\_match\_train.csv} & Supervised bundle-product links for training. \\
    \texttt{data/bundles\_product\_match\_test.csv} & Test template where \texttt{product\_asset\_id} must be predicted. \\
    \bottomrule
  \end{tabular}
  \caption{Repository datasets used by training and inference.}
  \label{tab:datasets}
\end{table}

\subsection{Preprocessing Pipeline}

Script \texttt{preprocess\_data.py} provides:
\begin{enumerate}
  \item Schema checks for required columns and referential integrity.
  \item Cleaning of product description labels, including exclusion of predefined cosmetic/perfume categories.
  \item Stratified train/validation split by number of labels per bundle.
  \item Concurrent image download with retry logic and forbidden URL logging.
  \item JSONL manifest generation and preprocessing statistics.
\end{enumerate}

The script produces:
\begin{itemize}
  \item \texttt{manifests/train\_manifest.jsonl}
  \item \texttt{manifests/val\_manifest.jsonl}
  \item \texttt{labels.json}, \texttt{label2idx.json}, and \texttt{stats.json}
\end{itemize}

\subsection{Offline Augmentation}

Script \texttt{offline\_augment.py} generates deterministic, seed-controlled augmented views for both bundles and products. It applies resized crops, color jitter, optional horizontal flips, Gaussian blur, optional cutout (bundles), and JPEG quality perturbation. Augmented images and manifests are saved as:
>>>>>>> e9d5d43845f1bcb17b1a4f5c33b124b3b6907bf4
\begin{itemize}
  \item \texttt{bundles\_aug/*.jpg}, \texttt{bundles\_aug\_manifest.jsonl}
  \item \texttt{products\_aug/*.jpg}, \texttt{products\_aug\_manifest.jsonl}
\end{itemize}

<<<<<<< HEAD
La motivacion principal es robustecer embeddings ante la restriccion de pocas vistas por producto.
=======
This design is aligned with the one-image-per-product constraint, where robust product representations require synthetic view expansion.
>>>>>>> e9d5d43845f1bcb17b1a4f5c33b124b3b6907bf4

\section{Metodologia (CRISP-DM)}

\subsection{Marco CRISP-DM aplicado}

La metodologia del proyecto sigue CRISP-DM y se adapta directamente a los scripts del repositorio (Tabla~\ref{tab:crisp-map}).

\begin{table}[h]
  \centering
  \small
  \begin{tabular}{p{0.27\linewidth} p{0.65\linewidth}}
    \toprule
    \textbf{Fase CRISP-DM} & \textbf{Aplicacion en goofytex} \\
    \midrule
    \textbf{Business Understanding} & Definicion del objetivo bundle$\rightarrow$product y del formato de entrega (maximo 15 predicciones por bundle). \\
    \textbf{Data Understanding} & Inspeccion de CSVs, mapeo de imagenes y analisis de metadata (scripts en \texttt{preprocess\_data.py} y \texttt{src/utils}). \\
    \textbf{Data Preparation} & Limpieza de etiquetas, manifests train/val, descarga/validacion de imagenes y augmentacion offline determinista. \\
    \textbf{Modeling} & Entrenamiento retrieval multimodal con OpenCLIP, pruebas de familias de embedding alternativas y deteccion de regiones con YOLO. \\
    \textbf{Evaluation} & Evaluacion local de ranking en validacion y chequeos funcionales del pipeline por script/variante. \\
    \textbf{Deployment} & Generacion de \texttt{submission.csv}, configuracion con Hydra y ejecucion reproducible por CLI. \\
    \bottomrule
  \end{tabular}
  \caption{Mapeo CRISP-DM a componentes reales del repositorio.}
  \label{tab:crisp-map}
\end{table}

\begin{figure}[h]
  \centering
  \small
  \begin{tabular}{c}
    \fbox{\parbox{0.92\linewidth}{\textbf{Business Understanding}\\Objetivo de retrieval, contrato de salida y restricciones de ranking}}\\[2pt]
    $\downarrow$\\[2pt]
    \fbox{\parbox{0.92\linewidth}{\textbf{Data Understanding + Preparation}\\CSVs, imagenes, manifiestos, limpieza, augmentacion y features auxiliares}}\\[2pt]
    $\downarrow$\\[2pt]
    \fbox{\parbox{0.92\linewidth}{\textbf{Modeling}\\Embeddings (QRLite/OpenCLIP/SigLIP), deteccion YOLO, retrieval y post-procesado}}\\[2pt]
    $\downarrow$\\[2pt]
    \fbox{\parbox{0.92\linewidth}{\textbf{Evaluation + Deployment}\\Validacion local, export de submission y iteracion controlada por configuracion Hydra}}
  \end{tabular}
  \caption{Diagrama de alto nivel del flujo CRISP-DM en el proyecto.}
  \label{fig:crisp-flow}
\end{figure}

\subsection{Modelos de embeddings probados}

Durante la iteracion metodologica se trabajaron varias familias de embedding visual/multimodal. Entre ellas, la linea CLIP/SigLIP~\cite{radford2021clip,zhai2023siglip} estructura el backend principal de representacion. La Tabla~\ref{tab:embedding-families} resume su papel.

\begin{table}[h]
  \centering
  \small
  \begin{tabular}{p{0.30\linewidth} p{0.22\linewidth} p{0.40\linewidth}}
    \toprule
    \textbf{Familia de modelo} & \textbf{Estado} & \textbf{Uso en el flujo} \\
    \midrule
    \textbf{QRLite} & Probado & Alternativa ligera para prototipos de embeddings y comparacion de coste/inferencia. \\
    \textbf{OpenCLIP Marqo Fashion SigLIP} & Principal & Backend actual de entrenamiento e inferencia multimodal (scripts \texttt{train}, \texttt{infer} y \texttt{new\_infer}). \\
    \textbf{SigLIP} & Probado & Familia contrastiva adicional usada en pruebas de embeddings y zero-shot de categorias. \\
    \textbf{Torchvision CNNs} (ResNet18/ResNet50/EfficientNet-B0) & Baseline & Encoders visuales preentrenados para baseline y comparativas de retrieval. \\
    \bottomrule
  \end{tabular}
  \caption{Familias de embedding consideradas en la metodologia.}
  \label{tab:embedding-families}
\end{table}

\subsection{Deteccion de prendas con YOLO}

Para capturar objetos dentro del bundle se integra deteccion por regiones. El repositorio usa por defecto \texttt{kesimeg/yolov8n-clothing-detection} mediante \texttt{ultralyticsplus}, y la metodologia contempla pruebas con varias configuraciones/checkpoints YOLO~\cite{jocher2023yolo} para ajustar granularidad de crops.

\begin{table}[h]
  \centering
  \small
  \begin{tabular}{p{0.32\linewidth} p{0.24\linewidth} p{0.36\linewidth}}
    \toprule
    \textbf{Componente YOLO} & \textbf{Implementacion} & \textbf{Funcion en pipeline} \\
    \midrule
    YOLO clothing (familia YOLOv8) & \texttt{src/detection.py} & Deteccion de cajas por bundle y cache JSON reutilizable. \\
    Variantes de umbral (conf/IoU) & Hydra config & Control de precision/cobertura de cajas antes del retrieval por crop. \\
    Fallback imagen completa & \texttt{src.infer*} & Garantiza inferencia cuando no hay detecciones validas. \\
    \bottomrule
  \end{tabular}
  \caption{Uso metodologico de YOLO en la fase de modelado/inferencia.}
  \label{tab:yolo-components}
\end{table}

\subsection{Pipeline tecnico end-to-end}

\begin{figure}[h]
  \centering
  \small
  \begin{tabular}{c}
    \fbox{\parbox{0.92\linewidth}{\textbf{1) Ingesta}\\CSVs + imagenes locales + metadata (seccion, genero, timestamps)}}\\[2pt]
    $\downarrow$\\[2pt]
    \fbox{\parbox{0.92\linewidth}{\textbf{2) Embeddings}\\Encoder principal (OpenCLIP Marqo SigLIP) o variantes de embedding probadas}}\\[2pt]
    $\downarrow$\\[2pt]
    \fbox{\parbox{0.92\linewidth}{\textbf{3) Deteccion y retrieval}\\YOLO por crop + similitud coseno + indice de productos (FAISS opcional)}}\\[2pt]
    $\downarrow$\\[2pt]
    \fbox{\parbox{0.92\linewidth}{\textbf{4) Post-procesado}\\Filtro de genero, diversidad por categoria, score threshold, TS rerank}}\\[2pt]
    $\downarrow$\\[2pt]
    \fbox{\parbox{0.92\linewidth}{\textbf{5) Salida}\\\texttt{submission.csv} con ranking por bundle}}
  \end{tabular}
  \caption{Pipeline bundle$\rightarrow$product implementado en los scripts de inferencia.}
  \label{fig:end2end-flow}
\end{figure}

\subsection{Recuperacion y scoring}

La similitud principal se calcula por coseno sobre embeddings normalizados:
\begin{equation}
\mathrm{sim}(\mathbf{q}, \mathbf{P}) = \mathbf{q}\mathbf{P}^{\top}
\end{equation}

En entrenamiento (\texttt{src/models/retrieval\_openclip.py}) se usa perdida contrastiva bidireccional bundle$\leftrightarrow$product, con scheduler coseno, warmup, AMP opcional y mineria de negativos duros configurable.

\subsection{Variantes de inferencia implementadas}

\begin{itemize}
  \item \texttt{src.infer}: pipeline principal (OpenCLIP + YOLO + filtros).
  \item \texttt{src.new\_infer}: retrieval per-crop con TTA y FAISS opcional.
  \item \texttt{src.infer\_phase1}: filtrado por seccion + zero-shot por categoria.
  \item \texttt{src.infer\_top1}: estrategia compacta top-1 por crop.
  \item \texttt{src.infer\_openclip}: script legacy por argparse.
\end{itemize}

\section{Implementation Details}

\subsection{Configuration and Entrypoints}

The project uses Hydra~\cite{yadan2019hydra} with typed dataclasses (\texttt{src/config.py}). Main entrypoints are:
\begin{itemize}
  \item \texttt{python -m src.train}: OpenCLIP training.
  \item \texttt{python -m src.infer}: pretrained baseline inference + local validation split.
  \item \texttt{python -m src.infer\_openclip}: inference from trained OpenCLIP checkpoint.
\end{itemize}

Core runtime parameters include batch size, epochs, learning rate, AMP, gradient accumulation, recall cutoffs, and optional detection settings (\texttt{bbox\_*}).

\subsection{Code Architecture}

\begin{table}[h]
  \centering
  \small
  \begin{tabular}{p{0.37\linewidth} p{0.55\linewidth}}
    \toprule
    \textbf{Module} & \textbf{Responsibility} \\
    \midrule
    \texttt{src/train.py} & Hydra bootstrap and trainer dispatch. \\
    \texttt{src/models/retrieval\_openclip.py} & Main OpenCLIP training/validation loop and checkpointing. \\
    \texttt{src/infer.py} & Baseline embedding extraction, retrieval, and submission export. \\
    \texttt{src/infer\_openclip.py} & OpenCLIP checkpoint inference for test bundles. \\
    \texttt{src/detection.py} & YOLO detector wrapper and box extraction utilities. \\
    \texttt{src/utils/retrieval.py} & Similarity and top-K retrieval helpers. \\
    \texttt{src/utils/metrics.py} & Hit@K and Recall@K computation. \\
    \texttt{preprocess\_data.py} & Dataset integrity checks, split, manifests, and stats. \\
    \texttt{offline\_augment.py} & Offline augmentation and augmentation manifests. \\
    \bottomrule
  \end{tabular}
  \caption{Primary modules and responsibilities.}
  \label{tab:code-map}
\end{table}

\subsection{Training Loop Highlights}

The OpenCLIP training implementation includes:
\begin{enumerate}
  \item Dynamic device resolution (CPU/CUDA) and optional multi-GPU DataParallel.
  \item Cosine learning-rate scheduler with 10\% warmup.
  \item Optional mixed precision via \texttt{torch.cuda.amp}.
  \item Periodic logging, per-epoch validation, JSONL metric appends, and checkpointing (\texttt{epoch\_X.pt}, \texttt{best.pt}).
\end{enumerate}

\subsection{Inference and Fallback Logic}

Both inference scripts generate one row per predicted product and cap outputs to 15 per bundle. If bundle embeddings are unavailable (e.g., missing image), popular products from train links are used as fallback.

\section{Experimental Protocol}

\subsection{Local Evaluation Strategy}

The repository provides two local validation pathways:
\begin{enumerate}
  \item Baseline mode (\texttt{src/infer.py}) samples validation bundles with \texttt{infer.val\_ratio} and reports \texttt{hit@K}/\texttt{recall@K}.
  \item Training mode (\texttt{src/models/retrieval\_openclip.py}) evaluates \texttt{recall@K} after each epoch, where $K=\texttt{params.recall\_k}$.
\end{enumerate}

Products are indexed once per validation pass, bundles are encoded, and top-K nearest neighbors are retrieved by cosine similarity.

\subsection{Output Artifacts}

Main artifacts generated by the training/inference workflow:
\begin{itemize}
  \item \texttt{metrics.jsonl}: epoch-level train loss and validation recall.
  \item \texttt{best.pt}: best checkpoint by recall metric.
  \item \texttt{test\_submission.csv}: final prediction file for competition upload.
  \item \texttt{val\_metrics.json}: baseline run summary and local metrics.
\end{itemize}

\subsection{Recommended Ablations}

To measure impact of major components, the following ablations are recommended:
\begin{itemize}
  \item Image-only products vs. image+text product embeddings.
  \item Bundle full-image encoding vs. detection-guided region aggregation.
  \item No offline augmentation vs. deterministic offline augmentation.
  \item Single-GPU vs. multi-GPU, with fixed seed and equal effective batch size.
\end{itemize}

\section{Reproducibility}

\subsection{Environment}

Base dependencies are listed in \texttt{requirements.txt}: \texttt{pandas}, \texttt{numpy}, \texttt{pillow}, \texttt{tqdm}, \texttt{requests}, \texttt{hydra-core}, \texttt{torch}, and \texttt{torchvision}. Additional runtime packages required by advanced features are:
\begin{itemize}
  \item \texttt{open\_clip\_torch} for the OpenCLIP backend.
  \item \texttt{ultralyticsplus} for bundle clothing detection.
\end{itemize}

\subsection{Determinism Controls}

The code explicitly sets random seeds (\texttt{random}, \texttt{numpy}, \texttt{torch}). Offline augmentation uses stable hash-based seeds per asset and augmentation index. This enables repeatable preprocessing and consistent training behavior under fixed configuration.

\subsection{Hydra Configuration}

Configuration is split into:
\begin{itemize}
  \item \texttt{config/config.yaml}: training and inference parameters.
  \item \texttt{config/files/file.yaml}: dataset/image/cache paths.
\end{itemize}

Hydra output directories store run-specific artifacts, making experiments traceable without manual path bookkeeping.

\subsection{Submission Contract}

The final CSV must contain:
\begin{itemize}
  \item \texttt{bundle\_asset\_id}
  \item \texttt{product\_asset\_id}
\end{itemize}
with up to 15 rows per bundle, sorted by model confidence.

\section{Limitaciones y Riesgos}

\begin{enumerate}
  \item \textbf{Supervision visual limitada por producto:} muchos IDs tienen pocas vistas, lo que complica robustez ante cambios de pose y oclusion.
  \item \textbf{Dependencia de modelos externos:} OpenCLIP y YOLO requieren pesos externos y versionado consistente de entorno.
  \item \textbf{Coste computacional:} deteccion por crops y rerankings incrementan latencia frente a pipelines full-image.
  \item \textbf{Dependencias opcionales no unificadas:} FAISS y Ultralytics se instalan aparte de \texttt{requirements.txt}.
  \item \textbf{Cobertura de tests limitada:} no hay una suite extensa de tests automatizados para todas las variantes de inferencia.
\end{enumerate}

\section{Conclusiones}

La documentacion de goofytex queda alineada con la estructura actual del repositorio, especialmente en rutas de utilidades y workflows de inferencia/reranking.

Las referencias de comandos, modulos y carpetas se han actualizado para reflejar la organizacion vigente en \texttt{src/}, \texttt{src/utils/}, \texttt{src/rerank/} y \texttt{documentacion/}.


\bibliographystyle{IEEEtran}
\bibliography{references}

\appendix
\section*{Apendice: Comandos Reproducibles}

\subsection{Descarga de Datos}
\begin{verbatim}
bash download.sh
\end{verbatim}

\subsection{Preprocesamiento}
\begin{verbatim}
python preprocess_data.py \
  --skip_download \
  --out_dir data/preprocessed \
  --val_ratio 0.1 \
  --seed 42
\end{verbatim}

\subsection{Aumento Offline}
\begin{verbatim}
python offline_augment.py \
  --bundles_manifest data/manifests/train_manifest.jsonl \
  --products_manifest data/product_dataset.csv \
  --products_images_dir data/product_images \
  --out_dir data/offline_aug \
  --bundles_num_augs 4 \
  --products_num_augs 2 \
  --img_size 224 \
  --seed 42 \
  --workers 8
\end{verbatim}

\subsection{Entrenamiento (Hydra)}
\begin{verbatim}
python -m src.train
\end{verbatim}

\subsection{Inferencia Baseline}
\begin{verbatim}
python -m src.infer
\end{verbatim}

\subsection{Inferencia OpenCLIP}
\begin{verbatim}
python -m src.infer_openclip \
  --checkpoint outputs/retrieval_openclip/best.pt \
  --submission-out outputs/retrieval_openclip/submission.csv
\end{verbatim}


\end{document}
